\section{Motivating Example}\label{sec:motivating-example}

This section presents a motivating feasibility case study that exposes a key reliability gap between \emph{verifier success} and \emph{developer-intended semantics} when synthesizing Move functions.
Our goal is not to claim statistical generality from this single case; rather, we use it to clarify what kind of feedback is actionable for Move-Prover-driven synthesis and why verifier pass alone can be misleading under partial specifications.

\subsection{Synthesis Task Setup}

We consider the synthesis of the \texttt{update\_performance\_statistics} function from the Aptos \texttt{stake} module.
Intuitively, this function updates validator performance counters after each block.
Concretely, it must update both:
(i)~counters associated with the block proposer (when the proposer index is present), and
(ii)~counters associated with validators that failed to propose, via an explicit loop over a vector of indices.

The task is constrained by a developer-provided formal specification that we treat as immutable.
The specification for \texttt{update\_performance\_statistics} (from \texttt{stake.spec.move}) contains a single \texttt{ensures} clause that constrains the successful-proposal counter:\footnote{We reproduce the specification verbatim from the Aptos Framework source tree (commit under review), eliding only unrelated pragmas.}
\begin{lstlisting}[language=Move,basicstyle=\small\ttfamily]
spec update_performance_statistics {
    requires chain_status::is_operating();
    aborts_if false;

    let validator_perf = global&lt;ValidatorPerformance&gt;(@aptos_framework);
    let post post_validator_perf = global&lt;ValidatorPerformance&gt;(@aptos_framework);
    let validator_len = len(validator_perf.validators);
    ensures (
        option::is_some(ghost_proposer_idx)
            && option::borrow(ghost_proposer_idx) < validator_len
    ) ==>
        (
            post_validator_perf.validators[
                option::borrow(ghost_proposer_idx)
            ].successful_proposals ==
            validator_perf.validators[
                option::borrow(ghost_proposer_idx)
            ].successful_proposals + 1
        );
}
\end{lstlisting}
Notably, the \texttt{ensures} clause constrains only the \texttt{successful\_proposals} counter; it does not contain a corresponding post-condition for \texttt{failed\_proposals}.
Because the specification is \emph{partial}---it captures safety properties and aborts-freedom but does not function as a full functional-equivalence oracle---an implementation that ignores \texttt{failed\_proposer\_indices} can still satisfy the written specification.
As we will show, this partiality produces ``verifier-passing but semantically incomplete'' implementations.

\subsection{Three Move-Prover Verification Idioms}\label{subsec:three-idioms}

When prompting a state-of-the-art coding LLM to synthesize a function body with only the signature, the fixed specification, and module context, the generated code often fails to pass the Move Prover.
Through inspection of prover traces, we observe that success requires satisfying three prover-specific idioms that are not native Move syntax and are sparsely reflected in public Move code:

\begin{enumerate}
    \item \textbf{Ghost-variable updates.}
    The specification references module-level ghost variables (e.g., \texttt{ghost\_valid\_perf} and \texttt{ghost\_proposer\_idx}).
    The body must synchronize them using inline \texttt{spec\,\{\,update\,\ldots\,\}} blocks:
    \begin{lstlisting}[language=Move,basicstyle=\small\ttfamily]
spec { update ghost_valid_perf = validator_perf; };
    \end{lstlisting}
    A regular \texttt{let} binding creates only a local Move variable and is invisible to the prover's specification state.

    \item \textbf{Loop invariants at the \texttt{while} header.}
    The Move Prover requires invariants to be embedded inside the \texttt{while}-header expression:
    \begin{lstlisting}[language=Move,basicstyle=\small\ttfamily]
while ({ spec { invariant ... }; condition }) { body }
    \end{lstlisting}
    Placing invariants inside the loop body triggers a compilation-level prover error.

    \item \textbf{Overflow assumptions.}
    Each \texttt{u64} increment must be preceded by an explicit assumption to keep the generated verification conditions tractable:
    \begin{lstlisting}[language=Move,basicstyle=\small\ttfamily]
spec { assume x + 1 <= MAX_U64; };
    \end{lstlisting}
    Without such assumptions, the prover emits unbounded arithmetic-overflow verification conditions and may exceed the global timeout.
\end{enumerate}

These idioms are difficult for general-purpose LLMs to infer from syntax alone: the surface-level code compiles, but the proof obligations fail because the body lacks the idiom-level structure the prover expects.

The only exception in our traces is GPT~5.5, which passes verification in zero-shot without producing any of the three idioms.
This is possible because its generated body is extremely simple: it updates only the successful-proposal counter and omits the entire \texttt{failed\_proposer\_indices} loop.
Without the loop, there are no loop invariants to place and no failed-proposal increments that would require overflow assumptions; without ghost-variable updates, the specification's \texttt{ensures} clause happens to be satisfiable because the omitted loop logic is the very part that would trigger the ghost-variable constraints.
Consequently, GPT~5.5 avoids the verification paths that require the idioms, but at the cost of producing a semantically incomplete body---an instance of the false-verified phenomenon we discuss in \S\ref{subsec:false-verified}.

\subsection{Why Raw Prover Feedback Is Not Actionable}

A natural attempt is to iterate: run the Move Prover on the generated body, feed the raw error trace back to the LLM, and ask for a revised implementation.
In practice, this ``generic feedback'' loop often fails to converge.
The prover outputs are machine-level (timeout messages, condition-generation errors, or counterexample traces) and do not directly map to the idiom-level transformations required by the Move Prover.

In the motivating case, generic feedback frequently triggers ``self-justification'' behavior: the model makes superficial syntax changes without discovering the missing ghost updates, invariant placement requirements, or overflow assumptions.
After one round of raw prover feedback, the auto-diagnoser misclassifies the root cause and suggests replacing \texttt{ghost\_proposer\_idx} in the \texttt{ensures} clause with the function parameter \texttt{proposer\_index}---whereas the actual fix is to add a \texttt{spec\,\{\,update\,\ldots\,\}} block in the body.
The error type shifts from timeout to ``loop invariants must be declared at the beginning of the loop header,'' but the loop still does not gain the correct invariant structure.

\subsection{Structured Diagnosis and the Emergent Second Failure Mode}

To make prover feedback actionable, MoVES employs a structured diagnosis step that maps raw failures to explicit repair instructions.
We illustrate its effect by running three representative LLMs on the same synthesis task under different feedback settings.

Table~\ref{tab:model-traces} summarizes the outcomes.
Columns denote the synthesis setting: \textbf{Zero-shot} = signature + spec only; \textbf{+Ctx} = zero-shot with additional module context; \textbf{+Diag.} = one round of structured diagnosis feedback.
Cell contents show the verifier result (timeout, error type, or PASS with prove time).

\begin{table*}[t]
\centering
\small
\begin{tabularx}{\textwidth}{lcccX}
\toprule
\textbf{Model} & \textbf{Zero-shot} & \textbf{+Ctx} & \textbf{+Diag.} & \textbf{Key Finding} \\
\midrule
Kimi & \xmark~timeout & \xmark~timeout & \cmark~PASS (77.8~s) & Diagnosis fixes ghost var, but body still semantically incomplete. \\
Claude Opus 4.7 & \xmark~timeout & \xmark~\texttt{old(..)} error & \cmark~PASS (72.7~s) & Structured diagnosis fixes all idioms in one round. \\
GPT 5.5 & \cmark~PASS (66.9~s) & --- & --- & \textbf{Code is semantically incomplete.} \\
\bottomrule
\end{tabularx}
\caption{Illustrative model traces on synthesizing \texttt{update\_performance\_statistics}.
This motivating example is qualitative ($N{=}1$ case) and intended to reveal failure modes rather than provide statistical evidence.}
\label{tab:model-traces}
\end{table*}

The concrete verifier timings for Kimi are as follows.
With signature and specification only (zero-shot, no module context), Kimi times out after 67.03~s; adding module context does not help (59.19~s, still timeout).
In both settings the prover stderr carries the compiler warning \texttt{Unused value of parameter failed\_proposer\_indices}.
After one round of structured diagnosis, Kimi's body passes in 77.79~s---yet the very same unused-parameter warning remains in the stderr, because the generated body still omits the loop that would consume \texttt{failed\_proposer\_indices}.

The most informative part is Kimi's structured-diagnosis run: while it succeeds in getting the Move Prover to accept the program, the generated code is still semantically incomplete.
The structured diagnosis identifies the ghost-variable idiom gap (enabling verifier pass), but the diagnostic taxonomy lacks a category for \emph{semantic-obligation omission}.
As a result, the model misinterprets the unused-parameter warning as something to be suppressed (e.g., renaming with a leading underscore), instead of inferring that the vector parameter must drive the missing loop logic.
This produces a second failure mode:
\begin{quote}
\em Structured diagnosis can improve verifier success, yet still fail to recover developer-intended semantics when the diagnosis does not capture semantic omissions.
\end{quote}

\subsection{The Semantic Gap: False-Verified Implementations}\label{subsec:false-verified}

A deeper inspection of the verifier-passing bodies reveals the same core issue: the implementations omit the loop over \texttt{failed\_proposer\_indices}.
Despite this semantic omission, the Move Prover accepts the generated code because the provided specification constrains only the aspects reflected by ghost-variable relations and certain safety properties, but does not fully enforce the intended behavior over the failed-proposer list.

Figure~\ref{fig:false-verified-code} contrasts Kimi's verifier-passing body (top) against the reference implementation (bottom).
The top version updates the proposer-related counters but never iterates through \texttt{failed\_proposer\_indices};
the bottom version includes the required loop, together with the three prover idioms from \S\ref{subsec:three-idioms}.

\begin{figure}[t]
\centering
\begin{lstlisting}[language=Move,basicstyle=\scriptsize\ttfamily]
/* Kimi + Struct. Diag.: verifier-passing but semantically incomplete */
spec { update ghost_proposer_idx = proposer_index; };
let validator_perf = borrow_global_mut<ValidatorPerformance>(@aptos_framework);
let validators = &mut validator_perf.validators;
if (option::is_some(&proposer_index)) {
    let idx = *option::borrow(&proposer_index);
    if (idx < vector::length(validators)) {
        let validator = vector::borrow_mut(validators, idx);
        validator.successful_proposals = validator.successful_proposals + 1;
    };
};
// NOTE: failed_proposer_indices is never used; no loop, no assume.

/* Reference / semantically complete obligation */
spec {
    update ghost_valid_perf = validator_perf;
    update ghost_proposer_idx = proposer_index;
};
let validator_perf = borrow_global_mut<ValidatorPerformance>(@aptos_framework);
let validator_len = vector::length(&validator_perf.validators);
if (option::is_some(&proposer_index)) {
    let idx = *option::borrow(&proposer_index);
    if (idx < validator_len) {
        let validator = vector::borrow_mut(&mut validator_perf.validators, idx);
        spec { assume validator.successful_proposals + 1 <= MAX_U64; };
        validator.successful_proposals = validator.successful_proposals + 1;
    };
};
let f = 0;
let f_len = vector::length(&failed_proposer_indices);
while ({
    spec { invariant f <= f_len;
           invariant forall j in 0..f:
               validator_perf.validators[
                   failed_proposer_indices[j]
               ].failed_proposals ==
               ghost_valid_perf.validators[
                   failed_proposer_indices[j]
               ].failed_proposals + 1; };
    f < f_len
}) {
    let validator_index = *vector::borrow(&failed_proposer_indices, f);
    if (validator_index < validator_len) {
        let validator = vector::borrow_mut(&mut validator_perf.validators, validator_index);
        spec { assume validator.failed_proposals + 1 <= MAX_U64; };
        validator.failed_proposals = validator.failed_proposals + 1;
    };
    f = f + 1;
};
\end{lstlisting}
\caption{Code comparison: Kimi's verifier-passing body (top) omits the \texttt{failed\_proposer\_indices} loop and the overflow assumptions present in the reference implementation (bottom).
`MAX\_U64' denotes the maximum value of type `u64', a constant recognized by the Move Prover.
The Move Prover accepts the top version because the partial specification does not constrain the failed-proposal update.}
\label{fig:false-verified-code}
\end{figure}

\subsection{Manual Semantic Coverage Audit (Illustrative)}\label{subsec:semantic-audit}

To separate \emph{verification} from \emph{semantic completeness}, we perform a manual semantic-coverage audit against five developer-oriented obligations:
(1)~updates successful-proposal counts,
(2)~handles the \texttt{failed\_proposer\_indices} loop,
(3)~updates ghost variables,
(4)~preserves the required \texttt{while}-invariant structure, and
(5)~adds overflow assumptions before arithmetic increments.

To isolate whether the bottleneck lies in code-generation capability or in diagnosis quality, we perform an oracle probe: we hand-write a structured diagnosis that explicitly names the three missing prover idioms and also flags the missing \texttt{failed\_proposer\_indices} loop as a semantic obligation.
Feeding this diagnosis into the same codegen prompt produces a verifier-passing body on the first attempt (76.95~s) that is semantically aligned with the reference implementation.
We stress that this hand-written diagnosis is used only as an upper-bound probe; it is not part of the automated MoVES pipeline.

Table~\ref{tab:semantic-coverage} summarizes the audit outcomes.

\begin{table*}[t]
\centering
\small
\begin{tabularx}{\textwidth}{lcccc}
\toprule
\textbf{Semantic Obligation} &
\textbf{GPT 5.5 Zero-shot} &
\textbf{Kimi + Struct.~Diag.} &
\textbf{Claude + Struct.~Diag.} &
\textbf{Human Oracle / Ref.} \\
\midrule
Updates successful-proposal counts & \cmark & \cmark & \cmark & \cmark \\
Handles \texttt{failed\_proposer\_indices} loop & \xmark & \xmark & \cmark & \cmark \\
Updates ghost variables & \xmark & \cmark & \cmark & \cmark \\
Preserves \texttt{while}-header invariant structure & \xmark & \xmark & \cmark & \cmark \\
Adds overflow assumptions before \texttt{+= 1} & \xmark & \xmark & \cmark & \cmark \\
\bottomrule
\end{tabularx}
\caption{Manual semantic coverage audit for the motivating example.
GPT 5.5 and Kimi both achieve verifier success while missing three or four semantic obligations.
The human-oracle diagnosis is an upper-bound probe and is not part of the automated pipeline.}
\label{tab:semantic-coverage}
\end{table*}

This audit reinforces two central takeaways:
\begin{enumerate}
    \item \textbf{Verifier pass is necessary but not sufficient} for semantic correctness under partial specifications.
    \item \textbf{Structured diagnosis must cover semantic-obligation recovery}, not only prover-facing idiom fixes.
    Otherwise, the generator may reach proof acceptance while still missing essential developer intent (as observed in Kimi and GPT 5.5).
\end{enumerate}

\subsection{What This Motivating Example Drives}

Driven by the above failure modes, we design MoVES around three requirements:
\begin{enumerate}
    \item a dual-role architecture that separates \emph{code generation} from \emph{error diagnosis},
    \item a Move-Prover-aware diagnosis taxonomy that translates prover failures into idiom-level repair instructions, and
    \item evaluation using semantic obligations beyond verifier acceptance to detect false-verified implementations.
\end{enumerate}

Together, these elements address the reliability gap highlighted by the motivating case: \emph{how to make verification feedback actionable without assuming that prover acceptance implies semantic completeness}.
